\documentclass[11pt]{article}
\usepackage[letterpaper,margin=0.85in,top=1.05in,bottom=0.7in]{geometry}
\usepackage{amsmath,amssymb}
\usepackage{xcolor}
\usepackage{fancyhdr}
\usepackage{enumitem}
\usepackage{lastpage}
%\usepackage{newtxtext,newtxmath}

\definecolor{accent}{HTML}{274C77}

\setlength{\parindent}{0pt}
\setlength{\parskip}{0pt}
\setlength{\headheight}{42pt}
\setlength{\headsep}{18pt}

\pagestyle{fancy}
\fancyhf{}
\fancyhead[L]{%
  \begin{minipage}[t]{0.56\textwidth}
  \textbf{CSCI 5526}\\[-1pt]
  Computational Tools for Multiscale Problems
  \end{minipage}}
\fancyhead[R]{%
  \begin{minipage}[t]{0.38\textwidth}\raggedleft
  August 20, 2026\\[-1pt]
  \textbf{Due:} 08/27/2026 2pm
  \end{minipage}}
\renewcommand{\headrulewidth}{0.6pt}
\renewcommand{\headrule}{\hbox to\headwidth{\color{accent}\leaders\hrule height \headrulewidth\hfill}}

% Number only; no "Problem x" heading.
\newcommand{\problem}[1]{\vspace{0.6em}\noindent{\large\bfseries #1.}\hspace{0.45em}}

% Uniform blank response area after every problem.
\newcommand{\answerspace}{\par\vspace{3.5in}}

\begin{document}

\begin{center}
  {\fontsize{25}{30}\selectfont\bfseries Problem Set 1}\par
  \vspace{0.35em}
  {\large\color{gray!75!black}Floating-Point Representation and Computer Arithmetic}
\end{center}
\vspace{0.75em}

\problem{1}
Consider a (slightly bizarre) variation of floating point representation, where the exponent $n$ is represented by two bits and the fractional part of the mantissa, $f$, is represented by seven bits. What are the largest and smallest positive machine numbers that can be represented using this system (not including $\pm\infty$ or $\pm 0$)? Give your answers as both binary and decimal numbers.
\answerspace

\problem{2}
Consider a variation of floating point representation in which the exponent is represented by $d$ bits and the fractional part of the mantissa is represented by four bits. If the largest and smallest positive machine numbers in this system are $248$ and $1.5625\times10^{-2}$, respectively, what is $d$?
\answerspace

\newpage

\problem{3}
Consider a computer that uses an arithmetic system that simply stores numbers in $n$-bit memory locations by dicing up the desired number range into $2^n-1$ even-size chunks. Explain why this is a Really Bad idea.
\answerspace

\problem{4}
Consider yet another variation of floating point representation with
\begin{itemize}[leftmargin=1.65em,itemsep=0.15em,topsep=0.35em]
  \item one sign bit ($1=$ negative),
  \item five exponent bits (without sign bits, but with the appropriate bias),
  \item and eight mantissa bits.
\end{itemize}
What is the decimal (base-10) value of the number that is represented as $1\ 10011\ 10001000$? Show your work.
\answerspace

\end{document}
